\section{The Spectral Sequence using Exact Couples}

\begin{definition}\label{def: Exact Couple giving A Morse Sps} Let $M$ be a smooth manifold with Morse data and
	\begin{equation*}
		\emptyset=:N_{-1}\subset N_0\subset N_1\subset \cdots \subset N_{m-1}\subset N_m=M,
	\end{equation*} be the filtration constructed in definition \ref{def:morseFiltr-aFiltrationOfSublevelsets}, i.e. $N_k=M^c$ where $c\in (k,k+1)$. Then for $p,q\in \N$ we define the Groups
	\begin{equation*}
		\begin{array}{llll}
			D^{p,q} & \coloneq K^{-(p+q)}(N_p)&
			~\text{and}~&D \coloneq \bigoplus_{p,q}D^{p,q}\\
			E^{p,q} & \coloneq K^{-(p+q)}(N_p,N_{p-1})
			&~\text{and}~ &E\coloneq \bigoplus_{p,q}E^{p,q}
		\end{array}
	\end{equation*}
	Now we have the (modified) exact sequence from the pair $(N_p,N_{p-1})$:
	\begin{center}

	\begin{adjustbox}{max width=\textwidth}
			\begin{tikzcd}
				\cdots\arrow[r] & K^{-(p+q-1)}(N_{p-1})\arrow[r,"b\circ \delta_{p,q}"]&K^{-(p+q)}(N_p,N_{p-1}) \arrow[r,"j_{p,q}^*"]&K^{-(p+q)}(N_p)\arrow[r,"i_{p,q}^*"]&K^{-(p+q)}(N_{p-1})\cdots
			\end{tikzcd}
	\end{adjustbox}		
	\end{center}
	(If $p+q=1$, we change the order from $b\circ \delta_{p,q}^*$ to $\delta_{p,q}^*\circ b$) This sequence is exact, as we can commute the isomorphism $\beta$ with $j_{p,q}$ to get the commutative diagram
	\begin{center}
		\begin{adjustbox}{max width=\textwidth}
			\begin{tikzcd}
				 K^{-(p+q-1)}(N_{p-1}) \arrow[dr,"b\circ \delta_{p,q}"]\arrow[r," \delta_{p,q}"]&K^{-(p+q)+2}(N_p,N_{p-1}) \arrow[r,"j_{p,q}^*"] \arrow[d,"b"] &K^{-(p+q)+2}(N_p)\arrow[d,"b"]\\
				 &K^{-(p+q)}(N_p,N_{p-1}) \arrow[r,"j_{p,q}^*"]&K^{-(p+q)}(N_p)
			\end{tikzcd}
		\end{adjustbox}
	\end{center}
	This is exact in the top row, as it is the exact sequence of a pair. Now its just a diagramm chasing to get the original statement: $y=b(x)\in \im(b\circ \delta_{p,q})$ is equivalent to $x\in \im(\delta_{p,q}^*)$ is equivalent to $j_{p,q}^*(x)=0$ which is equivalent to $b\circ j_{p,q}^*(x)=0$ which is equivalent to $j_{p,q}^*(y)=j_{p,q}^*\circ b(x)=0$. Hence, the kernal of $j_{p,q}^*$ equals the image of $b\circ \delta_{p,q}$. The same argument works in the case $q+p=1$
	With this we define the maps
	\begin{equation*}
		\begin{array}{lll|l}
			&&&\text{bidegree}\\
			\alpha&\coloneq \mathop{\oplus}\limits_{p,q}~i_{p,q}^*:&D\to D&(-1,1)\\
			\beta&\coloneq  \mathop{\oplus}\limits_{p,q}~\begin{cases}
				b\circ \delta_{p,q} & \text{ if }p+q\neq 1\\
				 \delta_{p,q} \circ b &\text{ if }p+q= 1
			\end{cases}:&D\to E&(1,0)\\
			\gamma&\coloneq \mathop{\oplus}\limits_{p,q}~j_{p,q}^*:&E\to D&(0,0)
		\end{array}
	\end{equation*}
	
	Obviously by definition the triangle is exact in every cornder:
	\begin{center}
		\begin{tikzcd}
			D\arrow[rr,"\alpha"]&&D\arrow[dl,"\beta"]\\
			&E \arrow[ul,"\gamma"]&
		\end{tikzcd}
	\end{center}
Finally we define $d\coloneq \beta\circ \gamma$.
\end{definition}

\begin{definition}[Categorie of Exact Couples]
	We define the categorie $\ExCoup$ as follows:
	The objects are exact couples, i.e. quintibles $(D,E,\alpha,\beta,\gamma)$ where $D,E$ are groups (or objects in an abelian categorie), such that the triangle is exact:
	\begin{center}
		\begin{tikzcd}
			D\arrow[rr,"\alpha"]&&D\arrow[dl,"\beta"]\\
			&E \arrow[ul,"\gamma"]&
		\end{tikzcd}
	\end{center}.Given another exact couple $(D',E',\alpha',\beta',\gamma')$ we define a morphism from one to another to be a touple of morphisms $(f:D\to D',g:E\to E')$ such that the following diagram commutes:
	\begin{center}
		\begin{tikzcd}
			D \arrow[r,"\alpha"]\arrow[d,"f"]& D \arrow[d,"f"] \arrow[r,"\beta"]&E \arrow[d,"g"] \arrow[r,"\gamma"]&D\arrow[d,"f"]\\
			D' \arrow[r,"\alpha'"]& D'  \arrow[r,"\beta'"]&E'  \arrow[r,"\gamma'"]&D'
		\end{tikzcd}
	\end{center}
\end{definition}
\begin{corollary}[An Endofunktor in the Categorie of Exact Couples]
	We can derive a new exact couple with the following definitions:
	\begin{equation*}
		\begin{array}{llc}
			\begin{array}{ll}
				E'&\coloneq \ker d \big/ \im d \\
				D'&\coloneq \alpha(A)
			\end{array}&
			\begin{array}{lll}
				\alpha'&:~ D'\to D'&a\mapsto \alpha(a)\\
				\beta' &:~ D'\to E'&a=\alpha(\tilde{a})\mapsto [\beta(\tilde{a})]\\
				\gamma' &:~ E'\to D' &[b]\mapsto\gamma(b) \\
			\end{array}
			&
			\begin{tikzcd}
				D'\arrow[rr,"\alpha'"]&&D'\arrow[dl,"\beta'"]\\
				&E' \arrow[ul,"\gamma'"]&
			\end{tikzcd}
		\end{array}
	\end{equation*} We again defnie $d'\coloneq \beta'\circ \gamma'$.
	The derifing of a new exact couple as defined above defines an endofunctor which er denote by $Der$.
\end{corollary}
\begin{proof}
	Let $(D,E,\alpha,\beta,\gamma)$ be the exact couple we started with and call the derived couple $(D',E',\alpha',\beta',\gamma')$.
	First we show the well definition of $\beta$. Let $a\in D'=\alpha(a_1)=\alpha(a_2)$. Then $a_1-a_2\in \ker\alpha=\im \gamma $. Hence $\beta(a_1)-\beta(a_2)=\beta(a_1-a_2)\in \im (\beta\circ \gamma)=\im d$. Furthermore, since $\beta(a_1)\in \im \beta=\ker\gamma \subseteq \ker d$ we have that the map $\beta'$ is well defined. For the well definition of $\gamma'$ we ssume that $[b_1]=[b_2]\in E'$. But then the difference of $b_1$ and $b_2$ is zero in $E'$ meaning in the image oif $d$, which is again contained in the image of $\beta$ and thereby in the kernal of $\alpha$. This menas that under $\alpha$ the image of $b_1$ and $b_2$ agree.
	Now to show that the derived couple is also exact one continuous the diagram yoga. We do exemplary the equality of the kernal of $\beta'$ with the image of $\alpha'$:
	Assume that $\beta'(x)=0$. That means that $\beta(\tilde{x})\in \im d$ with $\alpha(\tilde{x})=x$. We now want to show that $x\in \alpha^2D$, and do this by showing that $\tilde{x}=\alpha(z)$ for some $z$. Now since $\beta(\tilde{x})\in \im d$ we have that $\beta(\tilde{x})=d(y)=\beta\circ \gamma(x)$ for some $y$. In other words, $\tilde{x}-\gamma(y)\in \ker \beta=\im \alpha$. Hence, $\tilde{x}-\gamma(y)=\alpha(z)$ for some $z$. But now applying $\alpha$ to both side and using that $\alpha\circ \gamma=0$ we can deduce that $\alpha^2(z)=\alpha(\tilde{x})$. This shows that the kernal of $\beta'$ is contained in the image of $\alpha'$. The other direction of the inclusion is easier: Let $a=\alpha(\tilde{a}) \in D'$. Then $\beta'\circ \alpha'(z)=[\beta(z)]=\beta\circ \alpha (\tilde{a})=0$. Similary one checks for exactnes in the other corners.
	Now to show that we have an endofunktor we need to define what to do on morphisms. So assume that we have the two exact couples $(D,E,\alpha,\beta,\gamma)$ and $(\widetilde{D},\widetilde{E},\widetilde{\alpha},\widetilde{\beta},\widetilde{\gamma})$ with the morphisms $(f:D\to \widetilde{D},g:E\to \widetilde{E})$. We now need to define a morphism of exact couples $((f':D'\to \widetilde{D}',g:E'\to \widetilde{E}'))$ between the derived couples. For this we define:
	\begin{equation*}
		f':D'\to \widetilde{D}'\,,\,x\mapsto f(x) \quad \text{and}\quad g:E'\to \widetilde{E}'\,,\, [x]\mapsto [g(x)]
	\end{equation*} and check for well definition as maps and that the pair is a morphism of exact couples. The well definition of $f'$ follows directly from $f\circ \alpha=\widetilde{\alpha}\circ f$. For the well definition of $g'$ asume that $[x]=[y]\in E'$, meaning that $x-y\in \im d=\im \beta \circ \gamma$ hence for some $z$ we have that $x-y=\beta\circ \gamma(z)$ and hence $g(x-z)=g\circ \beta \circ \gamma (z)=\widetilde{\beta}\circ f\circ \gamma(z)=\widetilde{\beta}\circ \widetilde{\gamma}\circ g(z)\in \im \widetilde{d}$. This proofes the well definition.
\end{proof}

\begin{definition}[Spectral Sequence From an Exact Couple]
	We can now iterate this procedure to get a list of spaces with maps
	$ (E,d),(E',d'),(E'',d'')\dots$ or often written as $(E_1,d_1),(E_2,d_2),\dots$ which is called a \textbf{spectral sequence}. Notice how $E_{r+1}$ is given by $(E_r,d_r)$ whilst $d_{r+1}$ needs further information. Whilst the spectral sequence is defined inductively, the needed information is not captured in $(E_r,d_r)_r$ but rather in the exact couple. Or said differently, we get a sequence of exact couples $(E_r,D_r,\alpha_r,\beta_r,\gamma_r)$ from which we can get the spectral sequence. We can reformulate the maps as follows:
	\begin{equation*}
		D_r= \alpha^{r-1}(D),\quad E_r=\ker d_{r-1}\big/ \im d_{r-1},\quad  \alpha_r= \restr{\alpha}{ \alpha^{r-1}(D)} ,\quad \beta_r =[\beta \circ  \alpha^{1-r}], \quad \gamma_r= \gamma, \quad  d_r= \beta_r\circ \gamma_r
	\end{equation*} $\alpha^{1-r}$ means taking the preimage by alpha $1-r$ times. The well definition of this formulation was proven earlier. 
	The bidegree of the maps is easy to deduct from the formulations above:
	\begin{equation*}
		\begin{array}{lll}
			\alpha_r&\text{ has bidegree }& (-1,1)    \\
			\beta_r &\text{ has bidegree }& (-r+2,r-1)\\
			\gamma_r&\text{ has bidegree }& (0,0)     \\
			d_r     &\text{ has bidegree }& (r,+r+1)\\
		\end{array}
	\end{equation*}
\end{definition}
\begin{corollary}
	We can define the categorie of spectral sequences $\SpS$ where the objects consist of families of pairs $(E_r,d_r)_r\in \Z$ where $E_r$ are groups (or objects in an abelian categorie) and $d_r$ are endomorphisms in $E_r$, such that $E_{r+1}\cong \frac{\ker{d_r}}{\im d_r}$. Morphisms between two spectral sequences $(E_r,d_r)_{r\in \N}\to (\widetilde{E_r},\widetilde{d_r})_{r\in \N}$ are families of morphisms $(f_r:E_r\to \widetilde{E_r})$ such that for all $r$ we have that $f_r\circ d_r=\widetilde{d_r}\circ f_r$.
\end{corollary}
\begin{corollary}
	We have a functor from $\ExCoup$ to $\SpS$ given as follows: Let $(E,D,\alpha,\beta,\gamma)$ be an exact couple. Let $(E_r,d_r,\alpha_r,\beta_r,\gamma_r)=Der^r((E,D,\alpha,\beta,\gamma))$ where $Der^r$ denotes the r-th derivation of the exact couple. Then $(E_r,d_r\coloneq \beta_r\circ \gamma_r)_r$ is an exact sequence as described above. Assume that we have the two exact couples $(D,E,\alpha,\beta,\gamma)$ and $(\widetilde{D},\widetilde{E},\widetilde{\alpha},\widetilde{\beta},\widetilde{\gamma})$ with the morphisms $(f:D\to \widetilde{D},g:E\to \widetilde{E})$. Then the famliy $(Der^r(f))_r$ is a morphism of spectral sequences.
\end{corollary}
\begin{corollary}[Explicit Pages]\label{cor: explicit pages of exact sequences arising from exact couples}
	For a spectral seqquence arising from an exact couple defined as above we have a natural isomorphism
	\begin{equation*}
		E_{r+1}= \frac{\gamma^{-1}(\alpha^r D)}{\beta(\ker \alpha ^r)}
	\end{equation*}
\end{corollary}
\begin{proof}
	We proceed by induction starting with $n=1$.
	Here we have by exactnes:
	\begin{equation*}
		\ker(\beta \circ \gamma)=\gamma^{-1}(\ker(\beta))=\gamma^{-1}(\im(\alpha)),\quad \text{and}\quad
		\im(\beta\circ \gamma)=\beta(\ker(\alpha))\,.
	\end{equation*}
	Hence:
	\begin{equation*}
		E_2=\frac{\ker(d_1)}{\im(d_1)}=\frac{\ker(\beta\circ \gamma)}{\im(\beta \circ \gamma)}=\frac{\gamma^{-1}(\alpha D)}{\beta(\ker(\alpha))}\,.
	\end{equation*}
	Now Assume that
	\begin{equation*}
		E_r=\frac{\gamma^{-1}(\alpha^{r-1}D)}{\beta(\ker \alpha^{r-1})}
	\end{equation*}
	Then we have
	\begin{equation*}
		\ker{d_r}=\ker{\beta_r\gamma_r}=\gamma_r^{-1}(\im(\alpha_r))=\gamma_r^{-1}(\alpha^r(D))\,.
	\end{equation*} Now $\gamma_r:E_r\to D_r$ maps $[x]\mapsto \gamma(x)$, hence by induction assumption
	\begin{equation*}
		\ker{d_r}=\frac{\gamma{-1}(\alpha^nD)}{\beta(\ker \alpha)}\,.
	\end{equation*} Notice that $\alpha^{r-1}(\ker(\alpha_r))$ containes those elements in $D$, that are killed by the $\alpha^n$, in other words this equals $\ker \alpha^r$. Now $\beta_r=[\beta\circ \alpha^{r-1}]$ and hence using the induction assumptian again we can conclude:
	\begin{equation*}
		\im(\beta_r\gamma_r)=[\beta(\alpha^{r-1}\ker(\alpha_r))]=\frac{\beta(\ker\alpha^r D)}{\beta(\ker \alpha^{r-1})}\,.
	\end{equation*} Now using the isomorphism theorem we get the natural isomorphism from the statement.
\end{proof}



\begin{definition}
	In our case all pages of the spectral sequence inherit a double grading and the boundary maps too. Hence we will denote
	\begin{equation*}
		E_r=\bigoplus_{p,q}E_r^{p,q}\quad \text{and} \quad d_r=\oplus_{p,q}d_r^{p,q}
	\end{equation*}
\end{definition}


\begin{corollary}\label{cor: explicit pages of exact sequences arising from exact couples with grading}
	The grading carries through all pages and we can reformulate
	\begin{equation*}
		E^{p,q}_{r+1}=\frac{\gamma^{-1}\alpha^rD}{\beta(\ker \alpha^r)}\cap E^{p,q}_{r+1}
		=\frac{\gamma^{-1}\big((\alpha^rD)\cap D^{p,q}\big)}{\beta((\ker \alpha^r)\cap D^{p-1,q})}
		=\frac{\gamma^{-1}\alpha^rD^{p+r,q-r}}{\beta(\ker \restr{\alpha^r}{D^{p-1,q}})}
	\end{equation*}
\end{corollary}
\begin{corollary}[The Exact Sequence Stabilizes]
	For all $r\geq m+1$ it holds that
	\begin{equation*}
		E_{r}=E_{r+1}=\cdots
	\end{equation*} For this great enough $r$ we define $E_{\infty}\coloneq E_{r}$.
\end{corollary}
\begin{proof}
	Notice how our original $E^{p,q}$ is supportet in a strip regarding $p$ which is fancy languare for saying that $E^{p,q}=0$ for $p>0$ and for $p\geq m$. These conditinos carry through the whole spectral sequence, as suceeding pages are quotients of preceeding ones respecting the grading. Now the map $d_r$ has bidegree $(r,-r+1)$, and hence $d_r=0$ since the domain or codomain is zero , if the $p$ grading is greater then the $p-$length of the supporting strip meaning $p>m$ which gives us the claimes $r\geq m+1$.
\end{proof}
Now we refomulate the corollary \ref{cor: explicit pages of exact sequences arising from exact couples} for the infinity page. We have the constant sequence of groups $E_r=E_{r+1}=\cdots$ and by the corollary we have
\begin{equation*}
	\frac{\gamma{-1}(\alpha^rD)}{\beta(\ker{\alpha^r})}=\frac{\gamma{-1}(\alpha^{r+1}D)}{\beta(\ker{\alpha^{r+1}})}=\frac{\gamma{-1}(\alpha^{r+2}D)}{\beta(\ker{\alpha^{r+2}})}=\cdots
\end{equation*}
Hence in our later inspections of $E_{\infty}$ we can assume $r$ to be arbitrary big.
\begin{lemma}[The last page]
	In our spectral sequence, we have that
	\begin{equation*}
		E_{\infty}^{p,q}=\frac{\ker[K^{p+q}(M)\to K^{p+q}(N_{p-1})]}{\ker[K^{p+q}(M)\to K^{p+q}(N_{p1})]}
	\end{equation*}
\end{lemma}
\begin{proof}
	For each fixed $(p,q)$ there is an $r$ big enough such that $\alpha^r D^{p,q}=0$, since $\alpha^r$ has degree $-r$ in $p$, hence for $r>p$, $\alpha^r$ maps $D^{p,q}$ to something negative in $p$ and hence to the trivial group.
	Thus $D=\cup_{r\geq 0}\ker \alpha^r$. Now the denominator of $E^{p,q}_r$ equals $\beta(\ker( \restr{\alpha^r}{D^{p+q-1,q-r}})$ and hence for $r$ big enough this equals the image of $\restr{\beta}{D^{p-1,q}}$ which in turn is the kernal of $\restr{\gamma}{E^{p,q}}$. For the numerator, we start by looking at $(\alpha^r D)^{p,q}=\alpha^rD^{p+r,q-r}$. Now $D^{p+r,q-r}=K^{p+q}(N_{p+n})$. Hence for $n$ big enough (such that $N_{p+n}=M$) we can calculate $\alpha^n(K^{p+q}(N_{p+n}))=\im(K^{p+q}(M)\to K^{p+q}(N_p))$. This results in the description\todo{is das wahr mit gamma eingeschränkt auf Epq?}
	\begin{equation*}
		E^{p,q}_\infty=\frac{\gamma^{-1}\big(\im(K^{p+q}(M)\to K^{p+q}(N_p))\big)}{\ker\restr{\gamma}{E^{p,q}}}
	\end{equation*}
	Next we define the groups
	\begin{equation*}
		L^{p,q}\coloneq \im(K^{p+q}(M)\to K^{p+q}(N_p)) \cap \im \gamma= \im(K^{p+q}(M)\to K^{p+q}(N_p)) \cap \ker \alpha
	\end{equation*}
	And since we have for any $f$ and any subset $S$ of the  codomain of $f$ the equality $f^{-1}(S)=f^{-1}(S\cap \im f)$ holds,
	we have in our setting $ E^{p,q}_\infty=\frac{L^{p,q}}{\ker(\restr{\gamma}{E^{p,q}})}$. Now the kernal of the funktion
	\begin{equation*}
		\tilde{\gamma}:\gamma^{-1} (L^{p,q})\to L^{p,q};~x\mapsto \gamma(x)
	\end{equation*} equals $\ker \restr{\gamma}{E^{p,q}}$: Since $\gamma^{-1} (L^{p,q})\subseteq E^{p,q}$ the kernal of  $\tilde{\gamma}$ is clearly contained in $ker \restr{\gamma}{E^{p,q}}$. If on the other hand $\gamma(x)=0$ for $x\in E^{p,q}$ then $x\in \gamma^{-1}(0)$ and since zero is contained in $L^{p,q}$ we conclude that the two kernals are equal and hence we can apply the isomorphism theorem to get the isomorphism that is induced from $\gamma$:
	
	\begin{equation*}
		E^{p,q}_\infty=\frac{L^{p,q}}{\ker(\restr{\gamma}{E^{p,q}})} \to L^{p,q}\,.
	\end{equation*}
	Now finally we define the mapping
	\begin{equation*}
		\ker \left(K^{p+q}(M)\to K^{p+q}(N_{p-1})\right) \to L^{p,q}
		=\im(K^{p+q}(M)\to K^{p+q}(N_p)) \cap \ker \alpha\\
	\end{equation*} That maps an element $x\in K^{p,q}(M)$ to the pullback of the inclusion $N_p\to M$. This map is well defined, since applying $\alpha$ to the image of an element from the domain is the same as the pullback of $x$ via the inclusion $N_{p-1}\to M$ and equlals zero by definition of the domain. Hence the original map maps into the kernal of alpha which proofs the well definition. This line of reasoning also shows the surjectivity since every element in $K^{p,q}(N_p)$ that comes from the whole $M$ and maps to zero if restricted to $N_{p-1}$ can be expressed as the pullback from the whole $M$ restricted to those, that vanish over $N_{p-1}$. TIts easy to see that the kernal of this surjection equals the kernal of the pullback
	\begin{equation*}
		K^{p,q}(M)\to K^{p,q}(N_p)\,,
	\end{equation*} because ever element that pulls back to zero over $N_p$ also pulls back to zero over $N_{p-1}$, and the other containemend is trivially true.
	To the surjection we can again apply the isomorphism theorem to get the isomorphism
	\begin{equation*}
		\frac{\ker \left(K^{p+q}(M)\to K^{p+q}(N_{p-1})\right) }{\ker\left( K^{p,q}(M)\to K^{p,q}(N_p)\right)} \cong L^{p,q}\,,
	\end{equation*} that is induced by the pullback of the inclusion. Together with the $\gamma$ induced isomorphism from $E_\infty^{p,q}\to L^{p,q}$ we conclude the proof.
\end{proof}
\begin{definition}[Associated Graded Group]
	Let $G$ be a group together with a finite filtration of subgroups $\{F^i\}_{i=-1,\dots,-m}$ for $i=-1,\dots m$ such that 
	\begin{equation*}
		G=F^{-1}\subseteq F^0 \subseteq F^1\subseteq \cdots \subseteq F^{m-1}\subseteq F^m=0
	\end{equation*}
	Then we define the \textbf{associated graded group of $G$ with respect to the Filtration $\{F_i\}$} to be 
	\begin{equation*}
		\gr(G)\coloneq \bigoplus_{i=0}^{m} F^{i-1}\big/ F^{i}\, .
	\end{equation*} 
	From $\gr(G)$ we can rebuild $G$ by rebuilding $\{F^i\}$ up to extension problems in the following way:
	\begin{equation*}
		\begin{array}{ll}
			F^{m-1} & =F^{m-1}\big/F^m=\gr(G)_m \\
			F^{m-2} &  \text{ is an extension } 
			\begin{tikzcd}
				0 \arrow[r] & F^{m-1} \arrow[r] & F^{m-2} \arrow[r] & F^{m-2}\big/F^{m-1}=\gr(G)_{m-1} \arrow[r] & 0
				\end{tikzcd}\\
		F^{m-1}&  \text{ is an extension } 
		\begin{tikzcd}
			0 \arrow[r] & F^{m-2} \arrow[r] & F^{m-3} \arrow[r] & F^{m-3}\big/F^{m-2}=\gr(G)_{m-2} \arrow[r] & 0
		\end{tikzcd} \\
		\vdots & \cdots\\
		F^{-1}&  \text{ is an extension } 
		\begin{tikzcd}
			0 \arrow[r] & F^{0} \arrow[r] & F^{-1} \arrow[r] & F^{-1}\big/F^{0}=\gr(G)_{0} \arrow[r] & 0
		\end{tikzcd}
		\end{array}
	\end{equation*}
\end{definition}